\section{Methodology}

Our insight that modality constrains evaluation metrics drives our methodological design: capture modality via semantic embeddings, let clustering discover modality-driven families, then validate predictive power via historical train/test split.

\subsection{Feature Extraction Pipeline}

We extract four design features from each benchmark paper:

\begin{enumerate}
\item \textbf{Task Type:} Classification, generation, translation, question-answering (from Papers with Code taxonomy)
\item \textbf{Evaluation Metrics:} Primary metrics reported (accuracy, F1, BLEU, ROUGE, mAP, IoU, WER)
\item \textbf{Data Modality:} Image, text, audio, video, multimodal, tabular
\item \textbf{Dataset Size:} Number of examples in test/validation split
\end{enumerate}

\textbf{Extraction Protocol:} We developed a standardized protocol with objective decision rules (e.g., ``If paper reports BLEU, classify as text generation''). Two simulated annotators extract features independently from 20 benchmarks to measure inter-rater reliability via Cohen's kappa. Target: kappa $>0.80$ for scalability.

\textbf{Rationale for SentenceBERT over manual features:} While we extract categorical features for validation, we use SentenceBERT embeddings of benchmark descriptions for clustering. This captures semantic similarity enabling modality patterns to emerge from text rather than requiring pre-defined categories. Alternative approaches (one-hot encoding of manual features) lose semantic similarity; purely manual coding risks subjectivity (kappa $<0.80$).

\subsection{Coverage Family Discovery}

We cluster benchmarks using k-means on SentenceBERT embeddings:

\begin{enumerate}
\item Embed benchmark descriptions (title + abstract + task description) using SentenceBERT
\item Apply k-means clustering with k selected via silhouette score maximization
\item Measure intra-family similarity (average cosine similarity within clusters)
\item Validate cluster quality via modularity score on citation network
\end{enumerate}

\textbf{Historical Train/Test Split:} Pre-2023 benchmarks (2015--2022) form training set; 2023--2024 benchmarks form test set. This \textbf{breaks circular reasoning}---we test pure prediction rather than describing existing usage patterns. Alternative (cross-validation on full corpus) risks temporal contamination; random split loses temporal validation.

\textbf{Rationale for k-means:} Unsupervised discovery without labeled data. Silhouette score validates cluster quality. Hierarchical clustering provides dendrograms but k-means scales better to 100+ benchmarks. Supervised approaches require coverage family labels we're discovering.

\subsection{Citation Co-Occurrence Measurement}

For each coverage family, we measure citation overlap:

\begin{enumerate}
\item Extract methods citing benchmarks in each family (via Semantic Scholar API)
\item Calculate Jaccard similarity: $|\text{methods citing both benchmarks}| / |\text{methods citing either}|$
\item Compare intra-family overlap (same family) vs random baseline (different families)
\end{enumerate}

\textbf{Success Criterion:} If coverage families truly capture design constraints, methods using benchmarks from the same family should cite each other at $>70\%$ rate (Prediction P1), exceeding random baseline ($\leq 50\%$) with statistical significance ($p<0.05$).

\subsection{Validation Steps}

We test four assumptions:

\textbf{H-E1 (Citation Classification):} SciBERT classifier achieves $>85\%$ precision distinguishing validation claims from baseline mentions (enables validation evidence extraction)

\textbf{H-M1 (Feature Extraction Objectivity):} Cohen's kappa $>0.80$ for manual feature extraction (enables scaling to 100+ benchmarks)

\textbf{H-M2 (Meaningful Families):} Intra-family similarity $\geq 0.60$ (confirms families aren't arbitrary)

\textbf{H-M3 (Historical Prediction):} Citation overlap $>70\%$ with $p<0.05$ vs random (validates main claim)
